Image restoration results are reported as comparisons between methods, and the
arm that does nothing is usually left out because its number is uninteresting.
This paper reports three comparisons from one project in which that number was
recorded anyway, and shows that in each case it decided what the comparison
could conclude. At toy scale, two small networks were trained under a
matched-step budget and compared with each other; both finished below the
degraded input they were asked to repair, at \ToyConv{} and \ToyTrans{} decibels
against \ToyDegraded{}, so no comparison between them was available, and the
project withdrew its own recorded conclusion once that was noticed. Raising the
training budget \BudgetFold{} on the better of the two cut the training objective by a
factor of \ToyLossDrop{} and moved restored fidelity \ToyLongRegression{}
decibels the wrong way. At subset scale, a released diffusion restoration model
was run to completion over \SubsetN{} crops with reference images present, and
lost to bicubic upsampling on \LossesBicubic{} of them, by \GapBicubicAbs{}
decibels on average; the gap is \NoiseMargin{} times the largest per-crop shift
the same model shows between two of its own runs on the same crops, so it is not
run-to-run variation. No perceptual measure was computed for those crops, which
is exactly why the comparison cannot be read as a verdict on the model: what it
establishes is that a fidelity-only readout ranks an interpolation above it, and
that reporting the model's \MethodPSNR{} decibels alone would have hidden that.
The third case is an earlier attempt over the full set of \CensusTotal{} inputs,
where a memory ceiling dropped \CensusSkipped{} of them, and the \CensusRan{}
that survived had no reference images and were never scored; which inputs
survived was fixed by output size rather than by design, so the measured set was
a property of the hardware. One recommendation follows from all three: print the
trivial arm beside the method, and treat any comparison the trivial arm wins as
a statement about the measurement rather than about the method.
